% =====================================================================================
% DRAFT (rewritten 2026-09-11 after the triple check, docs/VERIFICATION_2026-09-11.md) --
% cross-simulator validation on RMDC26 (GULLS) for the A&C revision.
%
% Every number is a macro from paper/gulls_macros.tex or paper/paper_macros.tex, and both tables are
% \input from paper/outputs/, all generated by paper/make_gulls_macros.py (fail-closed) from the
% committed artifacts under validation/gulls/. No number in this file is typed.
%
% Decisions encoded here: the shipped weights stay; the gap-aware finite-source checkpoint ships
% alongside (binml-gapaware.pt) with its own operating point. Which of the two finite-source
% checkpoints (ESPLMag2-trained fspl5s_g08 or its smooth-magnification rerun fspl5s_espl_g08) is
% named "the sidecar" is settled in paper/REVISION.md section 1.5 row 14; the macros \bmlSidecar*
% follow that choice.
%
% Where each block goes: [A] abstract, replacing its last two sentences; [B] Sec. limits, replacing
% the "legacy Roman-like schedule" paragraph; [C] a new section after Sec. limits; [D] Discussion;
% [E] model card and data availability. Before merging: add % AUTO-GENERATED by make_gulls_macros.py -- do not edit.
\newcommand{\bmlThreshold}{0.904}
\newcommand{\bmlGapSensNone}{0.93}
\newcommand{\bmlGapSensTwo}{0.86}
\newcommand{\bmlGapSensN}{100}
\newcommand{\bmlGullsNcat}{377{,}999}
\newcommand{\bmlGullsNcatSingle}{299{,}466}
\newcommand{\bmlGullsNcatPlanet}{35{,}647}
\newcommand{\bmlGullsNcatPlanetBin}{42{,}886}
\newcommand{\bmlGullsAmpMedSingle}{0.063}
\newcommand{\bmlGullsAmpFracSingle}{41}
\newcommand{\bmlGullsSubdaySingle}{33}
\newcommand{\bmlGullsSubdayPlanet}{0.31}
\newcommand{\bmlGullsSubdayPlanetBin}{0.15}
\newcommand{\bmlGullsEligibleN}{100{,}935}
\newcommand{\bmlGullsScoredPct}{56.4}
\newcommand{\bmlGullsBetweenPct}{22.0}
\newcommand{\bmlGullsOutsidePct}{3.6}
\newcommand{\bmlGullsLowcadPct}{17.9}
\newcommand{\bmlGullsNoBaseN}{84}
\newcommand{\bmlGullsN}{56{,}975}
\newcommand{\bmlGullsNSingle}{33{,}353}
\newcommand{\bmlGullsNPlanet}{11{,}388}
\newcommand{\bmlGullsNPlanetBin}{12{,}234}
\newcommand{\bmlGullsPlanetFracW}{22}
\newcommand{\bmlGullsQmedPlanet}{1.4\times10^{-4}}
\newcommand{\bmlGullsQmedPlanetBin}{1.2\times10^{-4}}
\newcommand{\bmlGullsQksExp}{-18}
\newcommand{\bmlGullsBinSrcPct}{55}
\newcommand{\bmlGullsDetCut}{60}
\newcommand{\bmlGullsWidePass}{10}
\newcommand{\bmlGullsWideFrac}{63}
\newcommand{\bmlGullsResonantPass}{85}
\newcommand{\bmlGullsBrightSingle}{10}
\newcommand{\bmlGullsFsLowSingle}{21}
\newcommand{\bmlGullsBrightPlanet}{23}
\newcommand{\bmlGullsFsLowPlanet}{30}
\newcommand{\bmlGullsBrightPlanetBin}{21}
\newcommand{\bmlGullsFsLowPlanetBin}{35}
\newcommand{\bmlGullsRhoPnn}{0.29}
\newcommand{\bmlGullsNrhoGtOne}{43}
\newcommand{\bmlGullsFaBright}{11.8}
\newcommand{\bmlGullsFaInside}{4.1}
\newcommand{\bmlOursNonpsplQmed}{0.12}
\newcommand{\bmlOursNonpsplStellarPct}{82}
\newcommand{\bmlGullsSeasonDays}{70.7}
\newcommand{\bmlGullsPauseHoursLo}{43}
\newcommand{\bmlGullsPauseHoursHi}{44}
\newcommand{\bmlGullsGapLo}{109}
\newcommand{\bmlGullsGapHi}{120}
\newcommand{\bmlGullsObservedDays}{685}
\newcommand{\bmlGullsSpanDays}{1{,}717}
\newcommand{\bmlGullsPartialBinsPct}{34}
\newcommand{\bmlGullsBudgetLo}{2.0}
\newcommand{\bmlGullsBudgetMid}{5.2}
\newcommand{\bmlGullsBudgetHi}{11.7}
\newcommand{\bmlGullsFaShipped}{35.6}
\newcommand{\bmlGullsRecShipped}{0.214}
\newcommand{\bmlGullsRecBinShipped}{0.248}
\newcommand{\bmlGullsRecAtRecallShipped}{0.504}
\newcommand{\bmlGullsRhoLoShipped}{0.402}
\newcommand{\bmlGullsRhoHiShipped}{0.258}
\newcommand{\bmlGullsFaGapaware}{11.6}
\newcommand{\bmlGullsRecGapaware}{0.295}
\newcommand{\bmlGullsRecBinGapaware}{0.313}
\newcommand{\bmlGullsRecAtRecallGapaware}{0.457}
\newcommand{\bmlGullsRhoLoGapaware}{0.046}
\newcommand{\bmlGullsRhoHiGapaware}{0.674}
\newcommand{\bmlGullsFaFspl}{4.8}
\newcommand{\bmlGullsRecFspl}{0.390}
\newcommand{\bmlGullsRecBinFspl}{0.421}
\newcommand{\bmlGullsRecAtRecallFspl}{0.378}
\newcommand{\bmlGullsRhoLoFspl}{0.031}
\newcommand{\bmlGullsRhoHiFspl}{0.288}
\newcommand{\bmlGullsFaFsplOne}{5.2}
\newcommand{\bmlGullsRecFsplOne}{0.369}
\newcommand{\bmlGullsRecBinFsplOne}{0.399}
\newcommand{\bmlGullsRecAtRecallFsplOne}{0.370}
\newcommand{\bmlGullsRhoLoFsplOne}{0.032}
\newcommand{\bmlGullsRhoHiFsplOne}{0.333}
\newcommand{\bmlGullsFaFsplFive}{8.7}
\newcommand{\bmlGullsRecFsplFive}{0.344}
\newcommand{\bmlGullsRecBinFsplFive}{0.372}
\newcommand{\bmlGullsRecAtRecallFsplFive}{0.445}
\newcommand{\bmlGullsRhoLoFsplFive}{0.040}
\newcommand{\bmlGullsRhoHiFsplFive}{0.464}
\newcommand{\bmlGullsFaNoisy}{12.1}
\newcommand{\bmlGullsRecNoisy}{0.298}
\newcommand{\bmlGullsRecBinNoisy}{0.320}
\newcommand{\bmlGullsRecAtRecallNoisy}{0.475}
\newcommand{\bmlGullsRhoLoNoisy}{0.053}
\newcommand{\bmlGullsRhoHiNoisy}{0.513}
\newcommand{\bmlGullsFaOnset}{5.3}
\newcommand{\bmlGullsRecOnset}{0.375}
\newcommand{\bmlGullsRecBinOnset}{0.407}
\newcommand{\bmlGullsRecAtRecallOnset}{0.381}
\newcommand{\bmlGullsRhoLoOnset}{0.031}
\newcommand{\bmlGullsRhoHiOnset}{0.326}
\newcommand{\bmlGullsNcheckpoints}{6}
\newcommand{\bmlGullsBudgetSlack}{0.05}
\newcommand{\bmlGullsColourGainLo}{0.042}
\newcommand{\bmlGullsColourGainHi}{0.068}
\newcommand{\bmlGullsTruthN}{6{,}600}
\newcommand{\bmlGullsTruthNSingle}{1{,}600}
\newcommand{\bmlGullsTruthNPlanet}{2{,}500}
\newcommand{\bmlGullsEpochExcess}{11}
\newcommand{\bmlGullsUndetPlanet}{44}
\newcommand{\bmlGullsUndetPlanetBin}{46}
\newcommand{\bmlGullsDetPlanet}{56}
\newcommand{\bmlGullsDetPlanetBin}{54}
\newcommand{\bmlGullsUndetPlanetW}{41}
\newcommand{\bmlGullsFailChi}{26}
\newcommand{\bmlGullsFloorVeto}{19}
\newcommand{\bmlGullsDetRescaled}{55}
\newcommand{\bmlGullsSinglePsplRefitPct}{2.1}
\newcommand{\bmlGullsRecDetPlanet}{0.474}
\newcommand{\bmlGullsRecDetPlanetBin}{0.557}
\newcommand{\bmlGullsRecGenPlanet}{0.380}
\newcommand{\bmlGullsRecGenPlanetBin}{0.404}
\newcommand{\bmlGullsFlagUndet}{24}
\newcommand{\bmlGullsFlagFloorVeto}{40}
\newcommand{\bmlGullsSingleMisfitK}{10}
\newcommand{\bmlGullsSingleMisfitN}{20}
\newcommand{\bmlGullsSingleMisfitLo}{30}
\newcommand{\bmlGullsSingleMisfitHi}{70}
\newcommand{\bmlGullsFloorFlagThreeLo}{38}
\newcommand{\bmlGullsFloorFlagThreeHi}{47}
\newcommand{\bmlGullsFloorFlagOneLo}{16}
\newcommand{\bmlGullsFloorFlagOneHi}{20}
\newcommand{\bmlGullsRecChiLo}{0.17}
\newcommand{\bmlGullsRecChiHi}{0.87}
\newcommand{\bmlGullsColDetOne}{0.534}
\newcommand{\bmlGullsColDetThree}{0.539}
\newcommand{\bmlGullsColUndetOne}{0.170}
\newcommand{\bmlGullsColUndetThree}{0.270}
\newcommand{\bmlGullsPrecSample}{0.694}
\newcommand{\bmlGullsPrecScored}{0.618}
\newcommand{\bmlSidecarRecDetPlanet}{0.325}
\newcommand{\bmlSidecarRecDetPlanetBin}{0.398}
\newcommand{\bmlGullsFloorDetLo}{70}
\newcommand{\bmlGullsFloorDetHi}{40}
\newcommand{\bmlGullsSubdayN}{1{,}385}
\newcommand{\bmlGullsSubdayTe}{0.13}
\newcommand{\bmlGullsSubdayFa}{0.7}
\newcommand{\bmlGullsSubdayFaW}{3.1}
\newcommand{\bmlGullsSubdayPerLo}{60}
\newcommand{\bmlGullsSubdayPerHi}{93}
\newcommand{\bmlGullsNoiseBright}{0.95}
\newcommand{\bmlGullsNoiseFaint}{2.64}
\newcommand{\bmlGullsNoiseBkg}{7}
\newcommand{\bmlGullsNoiseGlobal}{1.41}
\newcommand{\bmlSidecarThr}{0.957}
\newcommand{\bmlSidecarThrLo}{0.942}
\newcommand{\bmlSidecarThrHi}{0.962}
\newcommand{\bmlSidecarFa}{1.6}
\newcommand{\bmlSidecarRecPlanet}{0.24}
\newcommand{\bmlSidecarRecPlanetBin}{0.27}
\newcommand{\bmlSidecarFaLo}{1.3}
\newcommand{\bmlSidecarFaHi}{2.4}
\newcommand{\bmlSidecarPoolPrev}{6.0}
\newcommand{\bmlSidecarPoolCompl}{0.677}
\newcommand{\bmlOccN}{6{,}024}
\newcommand{\bmlOccFaObs}{4.9}
\newcommand{\bmlOccFaOne}{3.6}
\newcommand{\bmlOccRecObs}{0.35}
\newcommand{\bmlOccRecOne}{0.26}
 to paper.tex, run
% make_gulls_macros.py from build.sh, hash the GULLS artifacts into paper/results/MANIFEST.json, and
% add the bibliography entries listed at the end.
% =====================================================================================

% ---------------------------------------------------------------- [A] abstract (replacement)
Validation on our own simulator uses a continuous 15-minute schedule and an adopted $0.02$~mag
detectability floor, so the numbers above are benchmark values rather than \Rom{} yield forecasts.
On RMDC26, an independent GULLS simulation of the survey ($\bmlGullsN$ selected events, none used
for training), the shipped checkpoint flags $\bmlGullsFaShipped\%$ of single lenses as anomalous,
because that simulation's F146 observing pauses are absent from our training seasons. Fine-tuning the
same network on our own simulations, first with gap augmentation and then with finite-source single
lenses, lowers this to $\bmlGullsFaFspl\%$ at the same threshold and raises planetary recall at a
$\bmlGullsBudgetMid\%$ single-lens false-alarm budget from $\bmlGullsRecShipped$ (shipped) and
$\bmlGullsRecGapaware$ (gap augmentation only) to $\bmlGullsRecFspl$; RMDC26 also guided that
diagnosis and the choice of checkpoint, so these values are optimistic for it. Applying our labelling
policy to RMDC26's noise-free curves, $\bmlGullsUndetPlanet$--$\bmlGullsUndetPlanetBin\%$ of its
selected planetary events carry no anomaly the policy would claim within one season, so we report
recall against both label sets. Code, weights, and reproducibility artifacts are released openly.

% ---------------------------------------------------------------- [B] Sec. limits, replacement paragraph
\paragraph{The schedule, not the sampling rate, is what the shipped model cannot survive.}
Our training grid is continuous within a season. The survey as simulated in the RMDC26 release is
not: F146 pauses for $\bmlGullsPauseHoursLo$--$\bmlGullsPauseHoursHi$~h per
$\bmlGullsSeasonDays$-day season, in six or seven pauses whose phases differ between its six
high-cadence seasons (only those near days 1, 35 and 69.5 recur), and the season is 1.3 days shorter
than our 72-day window. Essentially none of our training events contains an empty mid-season F146
bin, so the model's only prior for an empty token is the unrevealed future of a truncated season.
Imposing the measured pauses on our own held-out seasons, with nothing else changed, collapses the
shipped model (Table~\ref{tab:gap}): single-lens recall falls from $\bmlGapShippedPsplClean$ to
$\bmlGapShippedPsplSched$, mostly to \texttt{NonPSPL}, and \texttt{Flat} recall from
$\bmlGapShippedFlatClean$ to $\bmlGapShippedFlatSched$, mostly to \texttt{PeriodicVar}; the eruptive
and long-period variables and \texttt{NonPSPL} itself also degrade, and \texttt{NonPSPL} precision
falls from $\bmlGapShippedPrecClean$ to $\bmlGapShippedPrecSched$. A single gap of half an hour costs
nothing, but one gap of 1--2~h already lowers single-lens recall from $\bmlGapSensNone$ to
$\bmlGapSensTwo$ ($\bmlGapSensN$ events per class), the lost events going to \texttt{NonPSPL}. The remedy is
an augmentation, not a new model: blanking one to eight contiguous runs of 1--12~h in every band
during a 12-epoch warm-start fine-tune of the shipped weights raises held-out macro-F1 under the
measured schedule from $\bmlGapFoneShipped$ to $\bmlGapFoneRand$ with no clean-data cost
($\bmlGapFoneCleanShipped$ versus $\bmlGapFoneCleanRand$).
%% PENDING (schedule arms): one sentence on training with the measured per-season pauses instead of
%% random gaps (sched_seasons vs rand_norelabel, on our held-out and on RMDC26 by season), and one on
%% the relabel hazard: with the first season's pauses and the legacy onset grid,
%% $\bmlGapRelabelPct\%$ of NonPSPL-labelled training events were eligible for the caustic-in-gap
%% relabel, which is not truth-based (the recorded onset is when the anomaly first becomes
%% detectable, not where the caustic is); new schedule arms train with it off.
The shipped checkpoint's input contract is therefore \emph{continuous F146}; for a gapped schedule,
the gap-aware checkpoint of \S\ref{sec:gulls} is the one to use.

\begin{deluxetable*}{lcccccc}
\tablecaption{Our own held-out seasons ($\bmlGapHeldN$ scored events; a further 20\% of a
$\bmlGapPoolN$-event pool chooses each operating point) with the measured RMDC26 pauses imposed:
event $i$ receives the empty bins of high-cadence season $i \bmod 6$ in every band. ``Clean'' is the
same set without pauses. Labels are those of the unblanked season. All rows are 12-epoch warm-start
fine-tunes of the shipped weights on one natural-prior pool, except g08e12 (a different pool, the
starting point of the finite-source checkpoint in \S\ref{sec:gulls}); one seed each. Generated from
\code{validation/gulls/schedule\_finetune.json}.\label{tab:gap}}
\tablehead{\colhead{augmentation} & \colhead{AP clean} & \colhead{AP pauses} & \colhead{PSPL} &
           \colhead{Flat} & \colhead{NonPSPL} & \colhead{macro-F1}}
\startdata
shipped & 0.959 & 0.169 & 0.03 & 0.17 & 0.84 & 0.375 \\
gap augmentation, g08e12 (other pool) & 0.952 & 0.896 & 0.89 & 0.97 & 0.92 & 0.875 \\
random gaps & 0.951 & 0.910 & 0.91 & 0.98 & 0.92 & 0.896 \\
random gaps, relabel off & 0.952 & 0.911 & 0.91 & 0.97 & 0.92 & 0.899 \\
first-season mask, relabel off & 0.952 & 0.568 & 0.41 & 0.76 & 0.98 & 0.674 \\
measured seasons, relabel off & 0.948 & 0.921 & 0.91 & 0.97 & 0.94 & 0.899 \\

\enddata
\tablenotetext{a}{PSPL, Flat and NonPSPL are argmax recalls under the pauses. The first-season mask
arm trained on one season's pauses (the mask this work first took to be ``the'' schedule); it matches
one RMDC26 season in six.}
\end{deluxetable*}

% ---------------------------------------------------------------- [C] new section
\section{Cross-simulator validation on an independent \Rom{} simulation}\label{sec:gulls}

Every number so far comes from our own simulator. The Roman Galactic Exoplanet Survey Project
Infrastructure Team's RMDC26 release \citep{RGESPIT2026} is a simulation of the survey made with
GULLS \citep{Penny2013,Penny2019} by a different group, with its own implementation of the observing
schedule, photometry and Galactic population. We use it to test transfer: no RMDC26 event entered
training and no operating threshold was chosen on it. It was, however, used to diagnose the residual
false alarms, to set the upper end of the single-lens source-size prior, and to choose among
$\bmlGullsNcheckpoints$ fine-tuned checkpoints scored on it (one seed each), so the numbers for the
finite-source checkpoint are the best of several on this set and are optimistic.

\subsection{Data, selection, and input contract}

RMDC26 contains $\bmlGullsNcat$ events: $\bmlGullsNcatSingle$ single lenses (1S1L),
$\bmlGullsNcatPlanet$ planetary lenses (1S2L) and $\bmlGullsNcatPlanetBin$ planetary lenses with a
second source (2S2L). Every 1S2L and 2S2L event carries a planetary companion (median mass ratio
$\bmlGullsQmedPlanet$ and $\bmlGullsQmedPlanetBin$ on the scored events; similar but not identical
distributions, KS $p \sim 10^{\bmlGullsQksExp}$), the source is flagged binary in $\bmlGullsBinSrcPct\%$
of scored 2S2L events, and both planetary classes are pre-selected by GULLS' own detection statistic
($\bmlGullsDetCut$ or more for every such event). RMDC26 ships no pure binary-source class, so the
\citet{Gaudi1998} degeneracy of \S\ref{sec:sims} remains untested, and a \texttt{NonPSPL} call on a 2S2L
event may respond to its second source rather than its planet. RMDC26's anomalies are all planetary,
whereas $\bmlOursNonpsplStellarPct\%$ of our \texttt{NonPSPL} evaluation events are
stellar-mass-ratio binaries (median $q = \bmlOursNonpsplQmed$), so transfer recall is not comparable to
the in-distribution completeness of \S\ref{sec:results}.

Selection uses catalogue parameters, never the model's output, and is applied identically to the
three classes. We keep events whose host single-lens peak amplitude, computed from the catalogue impact
parameter and blend fraction, is at least $0.1$~mag (the median single-lens event is
$\bmlGullsAmpMedSingle$~mag and $\bmlGullsAmpFracSingle\%$ pass) and whose Einstein timescale lies in
our prior's support, $1 \le t_{\rm E} \le 300$~d ($\bmlGullsSubdaySingle\%$ of RMDC26 single lenses are
sub-day, against $\bmlGullsSubdayPlanet\%$ and $\bmlGullsSubdayPlanetBin\%$ of the planetary classes).
That leaves $\bmlGullsEligibleN$ events. Of these, $\bmlGullsScoredPct\%$ ($\bmlGullsN$:
$\bmlGullsNSingle$ 1S1L, $\bmlGullsNPlanet$ 1S2L, $\bmlGullsNPlanetBin$ 2S2L) peak inside one of the
six high-cadence seasons and are scored; $\bmlGullsBetweenPct\%$ peak in one of the
$\bmlGullsGapLo$--$\bmlGullsGapHi$-day gaps between seasons (\S\ref{sec:gulls:seasons}),
$\bmlGullsOutsidePct\%$ before or after the survey, $\bmlGullsLowcadPct\%$ in one of the four
low-cadence seasons (outside our support), and $\bmlGullsNoBaseN$ lack a usable baseline. The
amplitude cut is on the \emph{host} lens: it removes most wide-orbit planets, whose host is barely
magnified ($\bmlGullsWidePass\%$ of 1S2L events with $s > 2$ pass, against $\bmlGullsResonantPass\%$ at
$0.5 \le s < 2$), so the scored planetary events are dominated by resonant and central-caustic
geometries. Some scored inputs lie outside our training priors: the baseline is brighter than our
20~mag limit for $\bmlGullsBrightSingle\%$ of single lenses and $\bmlGullsBrightPlanet$--$\bmlGullsBrightPlanetBin\%$
of planetary events, and the blend fraction is below our prior's minimum for
$\bmlGullsFsLowSingle$--$\bmlGullsFsLowPlanetBin\%$ of events.

The model is given the 72-day window starting at the season boundary in all three bands. The
baseline is the median F146 magnitude more than $5\,t_{\rm E}$ from the peak over the whole survey,
using the catalogue $t_0$ and $t_{\rm E}$, because the catalogue baseline in this release is offset
from the observed quiescent flux; a real-time pipeline would not have either. Observations are pooled
to one value per 15-minute epoch before binning, as the training cache is built; binning by raw
observation instead moves the shipped checkpoint's single-lens false-alarm rate by about ten points,
so the operating point is sensitive to this choice. Two sampling differences survive preprocessing:
the pauses, and partial bin occupancy, since RMDC26's colour visits displace one of eight F146 epochs
in about $\bmlGullsPartialBinsPct\%$ of the 2-h bins, an occupancy our training data never contain.
Setting that occupancy to full on $\bmlOccN$ real RMDC26 inputs lowers the finite-source checkpoint's
single-lens false alarms from $\bmlOccFaObs\%$ to $\bmlOccFaOne\%$ and its planetary recall from
$\bmlOccRecObs$ to $\bmlOccRecOne$: the model reads partial occupancy as mild evidence of an anomaly.
(The pattern cannot be imposed faithfully on our binned simulations, because a bin's features would
still come from all its epochs; teaching it would need augmentation before binning.)

\subsection{Results}

Table~\ref{tab:gulls} gives the single-lens false-alarm rate at the frozen operating threshold and
planetary recall at matched single-lens false-alarm budgets, on the identical $\bmlGullsN$ events for
every checkpoint; the matched-budget comparison is the one a threshold shift cannot fake. The shipped
checkpoint flags $\bmlGullsFaShipped\%$ of single lenses. Every flagged event is, by construction,
called \texttt{NonPSPL}; across all single lenses the model's first choice splits between
\texttt{NonPSPL} and \texttt{PeriodicVar} with almost none called \texttt{PSPL}, the signature of the
empty mid-season tokens. Gap augmentation brings the false-alarm rate to $\bmlGullsFaGapaware\%$. Its
residual false alarms rise monotonically with $\rho/|u_0|$, the ratio of source radius to impact
parameter, from $\bmlGullsRhoLoGapaware$ for point-like sources to $\bmlGullsRhoHiGapaware$ above 3.
The high-ratio bins are high-magnification events with very small $|u_0|$, not large sources: only
$\bmlGullsNrhoGtOne$ scored single lenses have $\rho > 1$ (99th percentile $\bmlGullsRhoPnn$). The
likely cause is in our generator: its single lenses were point sources while its binaries sampled a
finite source, so a finite-source single lens, whose rounded peak resembles no training \texttt{PSPL},
was read as anomalous. Adding finite-source single lenses (a uniform, achromatic disc; $\rho$
log-uniform up to 5, binaries unchanged) and fine-tuning for a further 12 epochs from the gap-aware
checkpoint lowers the rate to $\bmlGullsFaFspl\%$ and the largest-ratio bin to $\bmlGullsRhoHiFspl$,
while planetary recall at the $\bmlGullsBudgetMid\%$ budget rises from $\bmlGullsRecGapaware$ to
$\bmlGullsRecFspl$.
%% PENDING (controls): the point-source control (same pool shape, recipe and 12 extra epochs) and the
%% smooth-magnification rerun -- how much of the gain is the physics, and whether the magnification
%% function's hand-off steps mattered.
Widening the \emph{binary} source-size prior at the same time was worse: false alarms
$\bmlGullsFaFsplFive\%$ instead of $\bmlGullsFaFspl\%$, recall $\bmlGullsRecFsplFive$ instead of
$\bmlGullsRecFspl$. The largest-ratio bins remain near $\bmlGullsRhoHiFspl$, so a large part of the
residual, not all of it, came from this missing physics. Among the scored single lenses, false alarms
are $\bmlGullsFaBright\%$ for baselines brighter than our 20~mag limit against $\bmlGullsFaInside\%$
inside it, so out-of-support brightness accounts for a further share.

Two other remedies were tested and failed. Matching GULLS' noise, whose F146 errors equal ours for
sources brighter than 20th magnitude and reach $\bmlGullsNoiseFaint$ times ours at 24--25~mag (a
$\bmlGullsNoiseBkg$-fold background variance reproduces this; a global factor does not), gave a
checkpoint that was worse at every matched budget and in false alarms ($\bmlGullsFaNoisy\%$; recall
$\bmlGullsRecNoisy$), although its recall at the frozen threshold was higher; with a single seed and a
single noise level we do not attribute a mechanism. Recording the anomaly onset on a finer grid (an
attempt made before the off-by-one error in its first implementation was found) moved recall at the
$\bmlGullsBudgetMid\%$ budget from $\bmlGullsRecFspl$ to $\bmlGullsRecOnset$.
%% PENDING (schedule arms): training on the measured per-season pauses instead of random gaps.
Two fine-tunes with the same recipe on different pools differ by $\bmlSeedSpreadLo$--$\bmlSeedSpreadHi$
in recall across the three budgets; differences of that size between single-seed checkpoints are not
resolved.

Removing the colour bands lowers planetary recall at matched budget by
$\bmlGullsColourGainLo$--$\bmlGullsColourGainHi$ for each of the three checkpoints ablated. The gain
does not come from the anomalies: on planetary events whose anomaly our policy can detect in F146,
colour changes recall at the $\bmlGullsBudgetMid\%$ budget only from $\bmlGullsColDetOne$ to
$\bmlGullsColDetThree$; on events without a detectable anomaly it raises the flag rate from
$\bmlGullsColUndetOne$ to $\bmlGullsColUndetThree$. We have not identified the mechanism (class-dependent
colours in RMDC26, or chromatic effects our F146 policy cannot see), so this is not evidence that the
colour bands carry anomaly information across simulators.

\begin{deluxetable*}{lcccccccc}
\tablecaption{Cross-simulator transfer on $\bmlGullsN$ RMDC26 events ($\bmlGullsNSingle$ 1S1L,
$\bmlGullsNPlanet$ 1S2L, $\bmlGullsNPlanetBin$ 2S2L), identical inputs for every row. ``FA'' is the
percentage of single lenses over the frozen threshold $\bmlThreshold$. ``Recall at'' columns give 1S2L
recall at the threshold that yields the stated single-lens false-alarm rate on this population (achieved
within $\bmlGullsBudgetSlack$ points); the budgets are the frozen-threshold false-alarm rates of three of
the checkpoints. The last four columns are the single-lens false-alarm rate at the frozen threshold in
bins of $\rho/|u_0|$. Fine-tunes are 12-epoch warm starts on our simulations, one seed each, each row
starting from the one above it except the control. Generated from
\code{validation/gulls/transfer\_tradeoff\_all.json}.\label{tab:gulls}}
\tablehead{\colhead{checkpoint} & \colhead{FA (\%)} & \colhead{recall at $\bmlGullsBudgetLo\%$} &
           \colhead{at $\bmlGullsBudgetMid\%$} & \colhead{at $\bmlGullsBudgetHi\%$} &
           \colhead{$\rho/|u_0| < 0.03$} & \colhead{0.3--1} & \colhead{1--3} & \colhead{$> 3$}}
\startdata
shipped & 35.6 & 38.3 & 0.122 & 0.214 & 0.313 & 0.144 \\
gap augmentation (g08e12) & 11.6 & 13.5 & 0.185 & 0.295 & 0.457 & 0.336 \\
\quad + point-source control & 6.4 & 8.6 & 0.244 & 0.361 & 0.485 & 0.374 \\
\quad + finite source & 4.8 & 7.6 & 0.265 & 0.390 & 0.533 & 0.383 \\
\quad + finite source, pauses (rec.) & 6.3 & 8.4 & 0.273 & 0.391 & 0.531 & 0.407 \\
\qquad same recipe, seed 2 & 4.7 & 5.8 & 0.232 & 0.349 & 0.503 & 0.386 \\
\qquad same recipe, seed 3 & 7.3 & 9.7 & 0.233 & 0.358 & 0.514 & 0.385 \\
% achieved single-lens false-alarm rate within 0.05 points of each budget

\enddata
\end{deluxetable*}

\paragraph{Operating point.} The frozen threshold was chosen on gap-free, point-source data. We chose
the finite-source checkpoint's threshold the way the paper chooses its own, at 90\% purity on our
simulated seasons (a finite-source held-out pool whose population-weighted \texttt{NonPSPL} prevalence,
$\bmlSidecarPoolPrev\%$, is close to our test set's $\bmlPriorPiZero\%$), with the measured RMDC26 pauses imposed and
without reading RMDC26. On the full pool this gives $\bmlSidecarThr$; the 20\% validation slices that
the paper's procedure would use scatter between $\bmlSidecarThrLo$ and $\bmlSidecarThrHi$ (68\%). At
$\bmlSidecarThr$ the checkpoint flags $\bmlSidecarFa\%$ of RMDC26 single lenses
($\bmlSidecarFaLo$--$\bmlSidecarFaHi\%$ across that scatter) with planetary recall
$\bmlSidecarRecPlanet$ / $\bmlSidecarRecPlanetBin$ (1S2L / 2S2L). The gap-aware checkpoint has no
usable operating point on a finite-source population: it reaches 90\% purity only on its few most
confident events.

\subsection{What ``recall'' means across two label ontologies}\label{sec:gulls:labels}

GULLS assigns classes by generator, subject to its own detection cut; BinML's classes are observational
(\S\ref{sec:detcond}). Because RMDC26 ships each event's noise-free light curve and per-epoch error, the
two can be reconciled: we applied our labelling rule, the $\Delta\chi^2$ thresholds and the $0.02$~mag
floor against the best-fitting static point-source single-lens model on the noise-free F146 curve, with
GULLS' errors as $\sigma$ and single lenses labelled exactly as in training, to the first
$\bmlGullsTruthNSingle$ scored single lenses and $\bmlGullsTruthNPlanet$ scored events of each planetary
class. Under that rule $\bmlGullsUndetPlanet\%$ of 1S2L and
$\bmlGullsUndetPlanetBin\%$ of 2S2L events carry no anomaly the policy would claim within one season:
$\bmlGullsFailChi\%$ of the planetary events fail the $\Delta\chi^2$ cut and $\bmlGullsFloorVeto\%$ pass
it but fall under the floor. Rescaling the $\Delta\chi^2$ cuts for RMDC26's $\bmlGullsEpochExcess\%$ more F146 epochs per
season moves these fractions by under a point. The figure depends on the floor, the single-season window
and our selection; it is not a statement about what a survey could detect. Recall against generator
labels mixes detectable and undetectable planets and is not bounded by the detectable fraction: the
finite-source checkpoint flags $\bmlGullsFlagUndet\%$ of the planetary events without a detectable
anomaly. On those with one, its recall at the frozen threshold is $\bmlGullsRecDetPlanet$ /
$\bmlGullsRecDetPlanetBin$ (1S2L / 2S2L) against $\bmlGullsRecGenPlanet$ / $\bmlGullsRecGenPlanetBin$ on
generator labels, and $\bmlSidecarRecDetPlanet$ / $\bmlSidecarRecDetPlanetBin$ at its own operating
point. Recall is limited by significance rather than amplitude: it rises from $\bmlGullsRecChiLo$ for
anomalies with $160 \le \Delta\chi^2 < 500$ to $\bmlGullsRecChiHi$ above $10^5$.

The rule also bears on the floor. The checkpoint flags $\bmlGullsFlagFloorVeto\%$ of the planetary events
whose anomaly is significant but below $0.02$~mag, at a similar rate however far below
($\bmlGullsFloorFlagThreeLo$--$\bmlGullsFloorFlagThreeHi\%$ across sub-floor amplitude bins); with F146
alone the rate falls to $\bmlGullsFloorFlagOneLo$--$\bmlGullsFloorFlagOneHi\%$, so much of it follows the
colour behaviour above. It also flags $\bmlGullsSingleMisfitK$ of the $\bmlGullsSingleMisfitN$ single
lenses whose static refit shows a significant sub-floor misfit (95\% interval
$\bmlGullsSingleMisfitLo$--$\bmlGullsSingleMisfitHi\%$; such misfits come from effects our static
point-source refit does not model, such as parallax): the response to sub-floor misfit is not specific
to planets. Lowering the floor from $0.02$ to
$0.005$~mag raises the fraction of 1S2L events with a claimable anomaly from $\bmlGullsDetPlanet\%$ to
$\bmlGullsFloorDetLo\%$; raising it to $0.05$~mag lowers it to $\bmlGullsFloorDetHi\%$. This is the
label side of the floor sensitivity; retraining under another floor is not covered here.

\subsection{Events whose peak falls between seasons}\label{sec:gulls:seasons}
%% PENDING (between-seasons v2): scorable fractions, wing-detectable fractions against in-season,
%% frozen-threshold and matched-budget recall with their thresholds, host-visible subset, what a
%% whole-mission model would add (events spanning seasons, the low-cadence seasons).

\subsection{The partial-season cascade on RMDC26}\label{sec:gulls:cascade}
%% PENDING (cascade rescan): half-day protocol, first-detectable onset recomputed on the noise-free
%% curve; premature rate, lag and detection against the in-house planetary strata (shipped model, our
%% simulator); single-lens alert burden; streaming purity at stated planetary prevalence (single-lens
%% contamination only; RMDC26 has no variable stars).

\subsection{What this test does and does not establish}

It establishes that the shipped model's failure on a gapped schedule is a property of the schedule,
reproducible without any GULLS data and remediable by augmentation, and that a specific missing piece of
physics in our generator accounted for a large part of the remaining single-lens false alarms. It does
not establish a \Rom{} planet yield: RMDC26's planetary population, detection cut and our amplitude
selection all differ from our training population; neither simulator contains real systematics; our
detectability policy uses a static point-source refit and a uniform source without limb darkening; the
fine-tunes' gap relabelling is approximate (\S\ref{sec:sims}); every fine-tune is a single seed; and
RMDC26 was used to choose among checkpoints. It does not test the binary-source degeneracy, for which
RMDC26 has no class. Sub-day single lenses, outside our timescale prior, are rarely flagged as anomalies
by the finite-source checkpoint ($\bmlGullsSubdayFa\%$ of $\bmlGullsSubdayN$ unweighted,
$\bmlGullsSubdayFaW\%$ weighted by event rate) but are read mostly as periodic variables
($\bmlGullsSubdayPerLo$--$\bmlGullsSubdayPerHi\%$ across checkpoints), not as microlensing; a model for
free-floating planets needs sub-day timescales in its prior, not a threshold change.

% ---------------------------------------------------------------- [D] Discussion paragraph
%% PENDING (between-seasons v2): the single-season input contract and the whole-mission model,
%% with the measured cost.

% ---------------------------------------------------------------- [E] model card / data availability
% Model card, "Input contract": continuous F146 within one 72-day season. For a gapped schedule such as
% RMDC26's, use binml-gapaware.pt at threshold \bmlSidecarThr (chosen on our own simulations with the
% measured RMDC26 pauses; 68% slice range \bmlSidecarThrLo-\bmlSidecarThrHi); the shipped threshold does
% not apply to it, and the frozen-threshold numbers above are for comparison only.
% Data availability, add: the RMDC26 release is on the Hugging Face Hub at RGES-PIT/MachineLearning,
% pinned to revision a338d5bab441b5caf551d2fea9469aadfdc81ec1; validation/gulls/rmdc26_scores.csv.gz
% holds every scored event's P(NonPSPL) for every checkpoint, from which the tables regenerate.

% ---------------------------------------------------------------- [F] bibliography to add before merging
% RGESPIT2026 (RMDC26 data release; citation and data-use terms to confirm with the RGES-PIT),
% Penny2013 (GULLS), Bozza2010 / Bozza2018 (VBBinaryLensing; already in refs.bib), STScI2025GBTDS (for
% the schedule; already in refs.bib).
